\subsection{A twenty-two-arrow category of type $\mathbb CP^2$}

\begin{theorem}\label{C:CP2-upper22}
There is a category with two objects and twenty-two arrows whose
classifying space is homotopy equivalent to $\mathbb CP^2$.
Its complete composition table and its comparison with the
twenty-three-element monoid in Theorem~\ref{B:CP23} are given in
\path{code/outputs/a_CP22_category.json}.
\end{theorem}
\begin{proof}
Use the explicit monoid $N$ from Theorem~\ref{B:CP23}.
Thus its minimal ideal is the band
\[
 K=I\times J,\qquad I=\{0,1,2\},\quad
 J=\{0,1,2,3,4\},
\]
and the remaining eight elements are $1,a,b,c,d,e,f,g$.
The row transformations of the seven nonidentity elements, in that
order, are
\[
 002,\quad002,\quad010,\quad202,\quad202,\quad202,\quad212.
\]
Every one preserves the two-element subset $I_0=\{0,2\}$.
Furthermore, a product of two elements outside $K$ that lies in
$K$ has its row in $I_0$.  Indeed none of the outside row
transformations is constant.  The only outside transformations
that take the value $1$ are $010$, $212$, and the identity, each
having singleton inverse image of $1$.  A composition of two
such nonconstant transformations cannot therefore be the constant
map with value $1$.

It follows that
\[
 P=(N\setminus K)\sqcup(I_0\times J)
\]
is an eighteen-element submonoid of $N$.  Products touching its
band stay in that band: left multiplication preserves $I_0$ and
right multiplication preserves the row.  Its minimal ideal is
exactly $I_0\times J$.  To see minimality directly, multiplication
on the two sides by this band reaches every one of its elements
from any fixed element of $P$.

Set $M=P^{\mathrm{op}}$ and let $S=I$.  Define a right $M$-action
by
\[
 i\cdot m=\sigma_m(i),
\]
where $\sigma_m$ is the original left row transformation in $N$.
The reversal of multiplication makes this a right action.
Every element of the minimal ideal of $M$ acts constantly on $S$.
The category $C=\mathcal C(M,S)$ is therefore explicitly specified
by
\[
 \End_C(a)=M,\qquad\End_C(b)=\{1_b\},\qquad
 \Hom_C(a,b)=I,\qquad\Hom_C(b,a)=\varnothing,
\]
with composition $i\circ m=\sigma_m(i)$.  Its number of arrows is
$18+1+3=22$.

Apply Theorem~\ref{C:constant-minimal-act-compression} to $(M,S)$.
The minimal ideal of $M$ has row set $J$ and column set $I_0$.
The compressed monoid has underlying set
\[
 (N\setminus K)\sqcup(J\times I).
\]
It is exactly $N^{\mathrm{op}}$: the identification is the identity
on $N\setminus K$ and sends the new element $(j,i)$ to the original
$x_{ij}$.  Products within the band, the two outside actions, and
outside products entering the band agree by their defining
formulas.  Consequently the actual comparison of classifying
spaces supplied by that theorem gives
\[
 BC\simeq B(N^{\mathrm{op}})\cong BN\simeq\mathbb CP^2.
\]
The last equivalence is the complete integral-resolution and
primitive-square result of Theorem~\ref{B:CP23}; thus this transfer
includes all homology degrees and the attaching-map information.

For a reproducible finite check, retain the source labels
$0,1,\ldots,5,11,12,\ldots,22$ in its numbering, take the opposite
table, and use the row action above.  The program
\path{code/a_CP22_category.py} checks all $972$ action equations,
all $6\,862$ composable associativity triples of the category,
and all $529$ product equations of the displayed compressed-monoid
isomorphism.  The same construction also works for the separately
certified source numbered $1651$; both literal models are retained
in the output file.  The separate verifier
\path{code/b_verify_CP22_category.py} imports no producer code and
independently reconstructs both submonoids, actions, categories,
and compressed multiplication tables; its results are recorded in
\path{code/outputs/b_CP22_category_independent.json}.
The additional independent reconstruction
\path{code/d_verify_CP22_category.py} also checks the comparison
functor on every composable pair.
\end{proof}
